% =====================================================================
\chapter{Funciones inyectivas, sobreyectivas y biyectivas}
\label{tema:3}
% =====================================================================

Este tema es el núcleo práctico del módulo: decidir si una función es inyectiva, sobreyectiva o biyectiva y \emph{demostrarlo}. Sigue la sección~12.2 de Hammack~\cite{hammack}, la sección~5.2 de Velleman~\cite{velleman} y las secciones~4.3 y~4.5 de~\cite{mcs}. Los ejercicios de clasificación siguen el problema~1 de la hoja~3 de Stanford CS103~\cite{cs103ps3}.

\section{Definiciones y método de demostración}

\begin{definicion}\label{def:t3-iny}
Sea $f\colon A\to B$.
\begin{enumerate}
  \item $f$ es \emph{inyectiva} si elementos distintos tienen imágenes distintas: para todos $a,a'\in A$, si $f(a)=f(a')$ entonces $a=a'$.
  \item $f$ es \emph{sobreyectiva} si todo elemento del codominio es imagen de alguno: para todo $b\in B$ existe $a\in A$ con $f(a)=b$; equivalentemente, $f(A)=B$.
  \item $f$ es \emph{biyectiva} si es inyectiva y sobreyectiva.
\end{enumerate}
\end{definicion}

En términos de fibras: $f$ es inyectiva si cada fibra $f^{-1}(\{b\})$ tiene \emph{a lo sumo} un elemento, sobreyectiva si tiene \emph{al menos} uno y biyectiva si tiene \emph{exactamente} uno. La \cref{fig:t3-flechas} muestra los cuatro casos posibles con diagramas de flechas.

\addtocounter{figure}{1}%

\paragraph{Cómo demostrar cada propiedad.}
\begin{itemize}[leftmargin=1.4em]
  \item \emph{Inyectividad}: se supone $f(a)=f(a')$ para $a,a'$ arbitrarios y se deduce $a=a'$ (demostración directa), o se supone $a\neq a'$ y se deduce $f(a)\neq f(a')$ (contrarrecíproco).
  \item \emph{No inyectividad}: basta un contraejemplo, $a\neq a'$ con $f(a)=f(a')$.
  \item \emph{Sobreyectividad}: se toma $b\in B$ arbitrario, se \emph{busca} un candidato $a$ (normalmente resolviendo la ecuación $f(a)=b$ en borrador), y se \emph{comprueba} que $a\in A$ y que $f(a)=b$.
  \item \emph{No sobreyectividad}: basta un $b\in B$ y una prueba de que la ecuación $f(a)=b$ no tiene solución en $A$.
\end{itemize}

\addtocounter{definicion}{1}%

\addtocounter{definicion}{1}%

\addtocounter{definicion}{1}%

\addtocounter{definicion}{1}%

\section{Criterios para funciones reales}

\begin{definicion}
$f\colon A\to\R$, con $A\subseteq\R$, es \emph{estrictamente creciente} si $x<x'$ implica $f(x)<f(x')$, y \emph{estrictamente decreciente} si $x<x'$ implica $f(x)>f(x')$. Es \emph{estrictamente monótona} si es una de las dos.
\end{definicion}

\begin{proposicion}\label{prop:t3-monotona}
Toda función estrictamente monótona es inyectiva.
\end{proposicion}

Si $f$ es derivable en un intervalo con $f'>0$ en todo punto, entonces es estrictamente creciente (consecuencia del teorema del valor medio, que se estudia en el módulo de cálculo de esta fase), y por tanto inyectiva. Para la sobreyectividad sobre un intervalo, el teorema de Bolzano (valores intermedios) permite concluir que una función continua toma todos los valores comprendidos entre sus límites; en este módulo se prefiere resolver $f(x)=y$ explícitamente.

\addtocounter{definicion}{1}%

\section{Funciones entre conjuntos finitos}

\begin{proposicion}\label{prop:t3-finitos}
Sean $A$ y $B$ finitos y $f\colon A\to B$.
\begin{enumerate}
  \item Si $f$ es inyectiva, $\abs{A}\le\abs{B}$; si es sobreyectiva, $\abs{A}\ge\abs{B}$.
  \item Si $\abs{A}=\abs{B}$, entonces $f$ es inyectiva si y solo si es sobreyectiva.
\end{enumerate}
\end{proposicion}

\begin{proposicion}\label{prop:t3-contar-iny}
Si $\abs{A}=k$ y $\abs{B}=n$, el número de funciones inyectivas $A\to B$ es $n(n-1)\cdots(n-k+1)$, que vale~0 si $k>n$. En particular, hay $n!$ biyecciones entre dos conjuntos de $n$ elementos.
\end{proposicion}

\addtocounter{definicion}{1}%

\addtocounter{definicion}{1}%

